\section{Introduction}

Quantum circuits are contractions of gate tensors and therefore tensor networks, and measurement probabilities are given by contractions.

\subsection{Quantum and Tensor Parallelism}

By its central axioms, quantum mechanics of multiple qubits is formulated by tensors capturing states and discrete time evolutions.
Motivated by the structural similarity, we investigate how quantum circuits can be utilized for the tensor-network based approach towards efficient and explainable AI in the \tnreason{} framework.

Potential Advantage: \emph{Quantum Parallelism} (see \cite[Section 3.2.5]{schuld_machine_2021}).
\begin{itemize}
    \item Evaluation of multiple function values by single circuit evaluation, we will relate it with the contraction of $\bencodingof{\exformula}$ here.
    \item Contraction perspective: Loop-tolerant efficient contractions.
    \item However: Need a generic encoding scheme to exploit this advantage, which is not known yet.
    \item We here only provide a scheme based on post-selection, which provides a quantum advantage only through amplitude amplification.
    Without amplitude amplification and post-selection of samples, the encoded distribution is always uniform.
    \item Further challenge: Application in NISQ devises \cite{preskill_quantum_2018}, instead of fault-tolerant quantum computers.
\end{itemize}
Comparison with classical side, which we can call \emph{Tensor Network Parallelism}:
Message-passing schemes for efficient contractions, but exact in limited cases.


